\begin{abstract}

\noindent\textbf{Background.}
Post-operative periorbital bruising is a frequent concern following eyelid surgery, with patient-visible sequelae that may influence recovery experience and satisfaction. Intraoperative irrigation with saline at different temperatures has been proposed as a simple means to reduce local bruising; however, rigorous evidence from controlled trials is lacking. This study aimed to evaluate the effect of irrigation saline temperature on post-operative periorbital bruising using a randomized, grader-blinded design.

\medskip
\noindent\textbf{Methods.}
Thirty-four patients undergoing eyelid surgery were enrolled; two were excluded (IDs~19 and~25), yielding 32 analysed patients (Condition~0: $n = 14$; Condition~1: $n = 18$). Post-operative bruising was assessed from standardised photographs by two independent, blinded graders across eight periorbital quadrants per patient (512 total observations). Outcomes were bruising depth and extent, each rated on an ordinal 0--4 scale. Inter-rater reliability was quantified by ICC(2,1) and quadratic-weighted Cohen's kappa with a pre-specified acceptability threshold of $\geq 0.75$. Treatment effects were estimated using a multi-model approach: linear mixed models (LMM), cumulative link mixed models (CLMM), and generalised estimating equations (GEE-Ordinal and GEE-Binary), all adjusted for procedure code and grader with a patient-level random intercept.

\medskip
\noindent\textbf{Results.}
No statistically significant treatment effect was detected for any outcome across any model. LMM estimates for the condition effect were: depth $\beta = -0.034$ (95\% CI: $-0.487$ to $0.420$, $p = 0.885$) and extent $\beta = 0.029$ (95\% CI: $-0.258$ to $0.316$, $p = 0.844$). CLMM estimates on the proportional odds scale were: depth OR $= 0.822$ ($p = 0.702$) and extent OR $= 1.006$ ($p = 0.987$). All effect sizes were negligible (Cohen's $d < 0.12$; Cliff's $\delta < 0.09$). Inter-rater reliability met the acceptability threshold for extent (ICC $= 0.888$) but not for depth (ICC $= 0.515$), with a systematic positive bias from Grader~1 (mean difference $+0.184$, $p < 0.001$). Post-hoc power was 5.5--7.2\%, indicating the trial was severely underpowered; minimum detectable effects were approximately 1.0~SD, requiring 571--2{,}786 patients per group to achieve 80\% power.

\medskip
\noindent\textbf{Conclusions.}
No meaningful difference in post-operative periorbital bruising was observed between treatment conditions across all statistical models and outcome definitions. The findings were consistent and robust to the analytic approach. However, the study was severely underpowered, and the absence of eye-level treatment assignment precluded recovery of the intended paired structure. Future trials should adopt eye-level randomisation with prospective recording of quadrant-to-eye mappings and pre-specified sample size calculations. The study remains blinded at the time of this analysis.

\end{abstract}
